\section{دو تمرین از چاپ تا کنترل قطعه}

هر دو مثال، طرح آموزشی همین جزوه‌اند. اندازه‌ها اسمی‌اند و برای کارگاه باید با چرم و ابزار واقعی بررسی شوند. در هر تمرین، ابتدا شابلون کاغذی را کنترل کنید و سپس روی دورریز برش آزمایشی بزنید.

\subsection{جاکارتی تک‌جیب روی یک برگ}

برگۀ تمرین در صفحۀ~\pageref{sheet:card} دو قطعه دارد. پشت \lr{A} با اندازۀ \lr{100 × 68}~میلی‌متر و جیب \lr{B} با اندازۀ \lr{100 × 52}~میلی‌متر. از هر قطعه یک عدد لازم است. دو ضلع عمودی و ضلع پایین به هم دوخته می‌شوند و لبۀ بالایی جیب دهانۀ ورود کارت است. مسیرهای خط‌چین حدود ۴~میلی‌متر از سه لبۀ برش فاصله دارند و محل تقریبی دوخت را نشان می‌دهند. فاصلۀ سوراخ‌ها را با ابزار دوخت خود تعیین کنید.

\begin{enumerate}
  \item برگه را در اندازۀ واقعی چاپ کنید. هر دو گیج \lr{H} و \lr{V} باید از مرکز نشانۀ ابتدا تا مرکز نشانۀ انتها ۱۰۰~میلی‌متر باشند. عرض دو قطعه و ارتفاع هرکدام را هم کنترل کنید.
  \item برای شروع، پشت \lr{1.2}~میلی‌متری و جیب \lr{0.8}~میلی‌متری را امتحان کنید. ضخامت اسمیِ درز پایین پیش از چسب و دوخت حدود \lr{2.0}~میلی‌متر است. برای پرداخت دهانۀ جیب، لبۀ بالایی آن را پیش از مونتاژ آماده کنید.
  \item پس از انتقال، هر قطعه را جداگانه ببرید. پایین و دو ضلع را روی هم بگذارید و خط دوخت را با اتصال واقعی چک کنید. اگر پشت قطعه را برای کم‌کردن برجستگی لویس می‌کنید، ابتدا همان دو لایه را روی دورریز امتحان کنید.
  \item با یک کارت واقعی، آزادی ورود، میزان بیرون‌ماندن کارت و دسترسی به آن را ببینید. فاصلۀ هندسیِ دو خط دوخت حدود ۹۲~میلی‌متر است. فضای مفید ممکن است با ضخامت، چسب و نخ کمتر شود. نتیجه را روی نمونۀ دوخته‌شده بسنجید.
\end{enumerate}

\subsection{پنل بزرگ کیف روی دو برگ}

این تمرین فقط پنل برش کیف را نشان می‌دهد. اندازۀ بیرونی آن \lr{280 × 140}~میلی‌متر است و دو برگۀ \lr{A4} آن در صفحات~\pageref{sheet:panel-left} و~\pageref{sheet:panel-right} آمده‌اند. دو برگ هرکدام ۱۵۰~میلی‌متر از عرض قطعه را نمایش می‌دهند و ۲۰~میلی‌متر هم‌پوشانی دارند. گوشه‌ها منحنی‌اند. جای اتصال پنل به بدنه یا آستر باید با الگوی محصول تعیین شود.

\begin{enumerate}
  \item هر دو برگ را در مقیاس ۱۰۰ درصد چاپ کنید. گیج افقیِ پایین هر برگ باید ۱۰۰~میلی‌متر باشد. فاصلۀ عمودی مرکز دو علامت تطبیق روی هر برگ ۱۵۲~میلی‌متر است.
  \item نوار سمت راست برگ چپ را در امتداد خط نقطه‌چین جدا کنید. لبۀ حاصل را روی خط نقطه‌چین برگ راست قرار دهید و مرکز دو علامت تطبیق را منطبق کنید. خط‌های پیوستۀ بالا و پایین باید در درز کاغذ بدون پله ادامه پیدا کنند. سپس برگ‌ها را به هم بچسبانید.
  \item عرض و ارتفاع شابلون مونتاژشده را اندازه بگیرید. برای برش چرمِ منعطف، شابلون را روی پوستِ بدون کشش ثابت کنید. برای چرم ضخیم، قبل از برش پنل، روی تکه‌ای هم‌ضخامت هم مسیر مستقیم و هم گردی گوشه را آزمایش کنید.
  \item ردِ اتصال دو برگ را به چرم منتقل نکنید. فقط محیط پیوسته را ببرید. اگر لبه به درز یا تا وارد می‌شود، محل و عرض لویس را با قطعه‌ای از همان ضخامت و با لایه‌بندی طرح تعیین کنید.
\end{enumerate}

\section{کنترل کیفیت پیش از تکرار}

\begin{table}[!htbp]
  \centering
  \caption{نشانه‌های خطا و اولین کنترل کارگاهی}
  \label{tab:checks}
  \begin{tabular}{@{}p{48mm}p{98mm}@{}}
    \toprule
    نشانه & اولین کنترل \\
    \midrule
    همۀ اندازه‌ها تقریباً به یک نسبت کوچک یا بزرگ‌اند & مقیاس چاپ و تنظیمات صفحه را از روی دو گیج کنترل کنید. چسباندن برگ با خطای مقیاس را متوقف کنید. \\
    خط پیوسته در درز کاغذ پله دارد & مرکز هر دو علامت تطبیق و هم‌راستایی بالا و پایین را هم‌زمان ببینید. \\
    قطعات جفت هم‌اندازه‌اند ولی در جهت اشتباه قرار دارند & سمت رو و پشت و الگوی چپ و راست را کنترل کنید. \\
    اندازه روی پشت و رو تفاوت دارد & تیغ ممکن است کج حرکت کرده باشد. زاویه، تیزی و سطح زیر قطعه را بررسی کنید. \\
    تاب‌برداشتن خط دوخت & عرض ناحیۀ لویس و ضخامت باقی‌مانده نزدیک سوراخ را بسنجید. \\
    قسمتی از قالب کامل نمی‌بُرد & سایش صفحه، سلامت تیغه، قرارگیری قالب و یکنواختی ضخامت را کنترل کنید. \\
    \bottomrule
  \end{tabular}
\end{table}

جدول~\ref{tab:checks} نقطۀ شروع عیب‌یابی است. برای تکرار، شابلون، جهت برش، ضخامت و نمونۀ پذیرفته‌شده را نگه دارید. پیش از سفارش قالب، واحد، ابعاد، بسته‌بودن مسیرها و تعداد قطعات را در فایل و نمونۀ کاغذی کنترل کنید.
